% ------------------------------------------------------------
% Page 102
% ------------------------------------------------------------

Je devais remplacer M. Galtier à l’Ecole d’Application à Mende, en 1955, à la veille de son
départ à la retraite. Les jours succèdent aux jours, dans une inquiétude grandissante, et les
problèmes d’installation de mon ménage sont loin d’être résolus.

\textbf{Visite de l’Inspecteur Primaire}

C’est alors que je suis assis à côté de mon petit élève Rémi que la porte de la sacristie s’ouvre,
livrant passage à l’Inspecteur Primaire, accompagné de mes collègues, l’instituteur de
Gatuzierès et celui de Fraissinet de Fourques. Je comprends que je vais subir les épreuves
orales du Certificat d’aptitude Pédagogique, le CAP.

Tout se passe bien. La critique pédagogique fort brève fut suivie d’un long entretien que je ne
pus éluder, tant il fut concret et déterminant, pour un processus d’installation dans une
situation bien particulière.

J’annonce mon mariage, et relate point par point mes démêlés avec un Maire qui ne veut rien
savoir et fuit ses responsabilités, et avec un Pasteur peu soucieux de ses devoirs de
propriétaire.
Je suis écouté sans être interrompu par mon Inspecteur, qui, à son tour et point par point,
analyse la situation :

« Vous n’aviez point à voir le Pasteur. C’est au Maire, usager ou locataire, qu’il appartient de
régler ses affaires avec le Pasteur ou la paroisse propriétaire. Enfin, vu l’urgence, je
comprends.
J’enregistre votre demande d’autorisation d’absence pour votre mariage, absence d’ailleurs
fort modeste. Le badigeon, vous ne pouvez pas l’obtenir en période scolaire. Cela se fait
pendant les grandes vacances. Pour ce qui est des WC le problème aurait dû être réglé depuis
longtemps, par la Mairie poussée par l’Administration, ce qui n’a peut-être pas été fait malgré
une insistance renouvelée. Je n’insisterai pas moi-même car c’est une école sans
avenir ! Allons voir cette fenêtre ! ».
J’indiquais que, après un double refus, j’étais décidé à la faire remplacer moi-même.

« ah ça non ! surtout pas ! je n’ai pas à en discuter avec le pasteur, mais avec le Maire. J’en
fais mon affaire ; dans moins d’une heure je vais le voir en repartant, et je vous promets que
vous l’aurez, cette fenêtre et tout de suite ! Et ce n’est pas vous qui la payerez, ni moi non
plus ! »

Quels furent ses arguments, je l’ignore, mais certainement très forts et convaincants. Dès le
lendemain un menuisier était là pour les mesures et, trois jours plus tard, la fenêtre était posée.

Avant son départ, l’Inspecteur me donna quelques conseils, dont je sus, par la suite, tirer
parti ainsi :
« J’ai vu du bois scié dehors et vous n’avez pas de bûcher. Rentrez donc ce bois, empilez-le
au fond de la classe, à l’intérieur de la porte cochère qu’on n’ouvre jamais. Vous aurez du bois
sec, et surtout vous diminuerez les courants d’air qui refroidissent cette grande salle. En un
mot vous êtes chez vous, et je suis votre chef direct ! »
